\chapter{Cronograma}
\label{cap:cronograma}

As atividades seguintes concentram-se no estudo sobre a composi\c{c}\~ao do conjunto de treinamento, contribui\c{c}\~ao central da disserta\c{c}\~ao, executado conforme o protocolo detalhado na Se\c{c}\~ao~\ref{sec:proposta-experimental-continuidade}, ``Proposta experimental para a continuidade da pesquisa'', e nas quest\~oes de avalia\c{c}\~ao deixadas em aberto nos Resultados.

A primeira etapa \'e ampliar e anotar o conjunto de dados: os 1.045 campos j\'a marcados no Google Earth, com uma imagem orbital cada, somados a novos ortomosaicos obtidos por drone, devem levar o acervo a mais de 1.200 imagens de campo, com o n\'umero de parcelas conhecido ap\'os a anota\c{c}\~ao. Nessa etapa tamb\'em ser\'a definido o protocolo de divis\~ao por grupos espaciais, que corrige a principal limita\c{c}\~ao da avalia\c{c}\~ao atual, com os grupos de teste separados antes de qualquer c\'alculo dos crit\'erios de sele\c{c}\~ao.

Com os dados organizados, ser\~ao executados os treinamentos previstos na proposta, os destinados ao c\'alculo da instabilidade e do F1 fora da amostra e os das quatro condi\c{c}\~oes comparadas, com as diferen\c{c}as estimadas por reamostragem e analisadas separadamente para imagens orbitais e de drone e por faixa de densidade de parcelas.

Paralelamente, ser\~ao examinadas duas quest\~oes. A primeira \'e o efeito da dimens\~ao de entrada, com a reavalia\c{c}\~ao dos checkpoints treinados com 1.280 pixels nessa mesma dimens\~ao. A segunda \'e a corre\c{c}\~ao dos identificadores atribu\'idos pela numera\c{c}\~ao: como o campo P11 disp\~oe de croqui digitalizado, com identifica\c{c}\~ao, gen\'otipo e bloco de cada parcela, os identificadores propagados entre imagens poder\~ao ser comparados \`a identifica\c{c}\~ao registrada no ensaio. Essa compara\c{c}\~ao exige, antes, estabelecer a correspond\^encia entre as 876 parcelas do croqui e as 217 da refer\^encia numerada utilizada.

A etapa final compreende a an\'alise dos resultados, a reda\c{c}\~ao da disserta\c{c}\~ao e a prepara\c{c}\~ao da defesa, conforme a Tabela~\ref{tab:cronograma-defesa}.

\begin{table}[htb]
\centering
\caption{Cronograma das atividades at\'e a defesa}
\label{tab:cronograma-defesa}
\small
\begin{tabular}{p{3.4cm}p{8.4cm}p{1.8cm}}
\hline
\textbf{Per\'iodo} & \textbf{Atividade} & \textbf{Situa\c{c}\~ao} \\
\hline
Setembro - outubro/2026 & Amplia\c{c}\~ao e anota\c{c}\~ao do acervo; confer\^encia e corre\c{c}\~ao das parcelas anotadas; defini\c{c}\~ao do protocolo de divis\~ao por grupos espaciais & Prevista \\
Outubro/2026 & Reavalia\c{c}\~ao dos \textit{checkpoints} treinados com 1.280 \textit{pixels} nessa mesma dimens\~ao de entrada; verifica\c{c}\~ao da composi\c{c}\~ao dos descritores de dados dos \textit{checkpoints} preservados & Prevista \\
Outubro - novembro/2026 & Treinamentos destinados ao c\'alculo da instabilidade e do F1 fora da amostra & Prevista \\
Novembro - dezembro/2026 & Treinamentos das condi\c{c}\~oes comparadas e avalia\c{c}\~ao nos grupos de teste & Prevista \\
Dezembro/2026 & Compara\c{c}\~ao dos identificadores propagados no campo P11 com o croqui digitalizado & Prevista \\
Janeiro/2027 & An\'alise dos resultados, estimativa das diferen\c{c}as por reamostragem e reda\c{c}\~ao dos cap\'itulos de resultados e discuss\~ao & Prevista \\
Fevereiro - mar\c{c}o/2027 & Revis\~ao do texto, prepara\c{c}\~ao da vers\~ao final e defesa da disserta\c{c}\~ao & Prevista \\
\hline
\end{tabular}
\end{table}

O principal risco do cronograma est\'a na etapa de treinamento: s\~ao 22 execu\c{c}\~oes de 120 \'epocas, 10 para o c\'alculo dos escores e 12 para as condi\c{c}\~oes comparadas, sobre um acervo que deve ultrapassar 1.200 imagens. Caso o tempo de processamento dispon\'ivel n\~ao comporte esse total, a redu\c{c}\~ao ser\'a feita primeiro na segunda repeti\c{c}\~ao da valida\c{c}\~ao cruzada por grupos, que reduz o n\'umero de treinamentos destinados ao c\'alculo dos escores de 10 para 5, e s\'o depois no n\'umero de sementes por condi\c{c}\~ao, preservando o contraste principal entre a sele\c{c}\~ao por instabilidade e os subconjuntos aleat\'orios de mesmo tamanho.
